	% ---
	% LISTA DE FIGURAS 
  %%%% Não é necessário alterar essa área, pois os dados são coletados automaticamente na adição de figuras 
  %%%% há um exemplo em 104.revisaoliteratura.tex
	% ---
	\renewcommand{\listfigurename}{Lista de figuras}
	\pdfbookmark[0]{\listfigurename}{lof}
	\listoffigures*
	\cleardoublepage
  % ---  
  
	% ---
	% LISTA DE QUADROS
	% ---
	\pdfbookmark[0]{\listofquadrosname}{loq}
	\listofquadros*
	\cleardoublepage
	% ---
  
	% ---
	% LISTA DE TABELAS
  %%%% Não é necessário alterar essa área, pois os dados são coletados automaticamente na adição de tabelas.
  %%%% há um exemplo em 105.materialemetodos.tex
	% ---
	\pdfbookmark[0]{\listtablename}{lot}
	\listoftables*
	\cleardoublepage
  % ---


  % ---
	% LISTA DE CÓDIGOS FONTE
  %%%% Não é necessário alterar essa área, pois os dados são coletados automaticamente na adição de códigos fonte
  %%%% há um exemplo em 105.materialemetodos.tex
	% ---
  \lstlistoflistings*
  
  
	% ---
	% LISTA DE ABREVIATUAS E SIGLAS
  %%%% Aqui é necessário alterar de acordo com as siglas utilizadas no seu projeto. Abaixo, alguns exemplos
	% ---
  \cleardoublepage
  \begin{siglas}
    \item[ABNT] Associação Brasileira de Normas Técnicas
    \item[UEPG] Universidade Estadual de Ponta Grossa
    \item[LCAD] Laboratório de Computação de Alto Desempenho
   \end{siglas}  
	% ---
	
	% ---
	% LISTA DE SÍMBOLOS
  %%%% Aqui é necessário alterar de acordo com símbolos no seu projeto. Abaixo, alguns exemplos
	% ---
  \begin{simbolos}
    \item[$ \Gamma $] Letra grega Gama
    \item[$ \Lambda $] Lambda
    \item[$ \zeta $] Letra grega minúscula zeta
    \item[$ \in $] Pertence
  \end{simbolos}
	% ---
  
  % ---
	% SUMÁRIO
  %%%% Não é necessário alterar essa área, pois os dados são coletados automaticamente de chapters, sections e subsections 
	% ---
	\pdfbookmark[0]{\contentsname}{toc}
	\tableofcontents*
	\cleardoublepage
	% ----------------------------------------------------------
	\textual
  % ---